% ============================================================================
% 2.2 堆内存分配器
% ============================================================================
\subsection{设计思路}

我们需要一个比 \texttt{.space} 伪指令更灵活的内存分配方式——
可以按需分配、用完释放、释放的空间再利用。

我们的实现是一个简化版的 K\&R 风格分配器：

\begin{enumerate}
    \item 在数据段预留一块 4096 字节的连续内存作为堆
    \item 每个内存块有一个 8 字节的头：前 4 字节存大小，后 4 字节存是否在用
    \item 分配时从堆头开始扫描，找第一个足够大的空闲块
    \item 释放时把块标记为空闲，并与左右相邻空闲块合并（coalescing）
\end{enumerate}

一个块的内存布局：

\begin{center}
\begin{tabular}{|c|c|c|}
\hline
size（4 字节） & used（4 字节） & 实际数据区（size 字节） \\
\hline
\multicolumn{2}{|c|}{8 字节头} & \\
\hline
\end{tabular}
\end{center}

整个堆的结构（假设分配了两块，中间一块空闲）：

\begin{center}
\begin{tabular}{|c|c|c|c|c|c|c|}
\hline
size 20 & used 1 & data 20B & size 100 & used 0 & data 100B（空闲） & data 40B \\
\hline
\end{tabular}
\end{center}

\subsection{heap 区的定义与初始化}

在 \texttt{hello.s} 的 \texttt{.data} 段里，我们这样定义堆：

\begin{lstlisting}[style=asm,caption=hello.s 中的 heap 定义]
.include "src/memory.inc"    ; 定义 HEAP_SIZE = 4096
.data
.align 4

.global heap
.global heap_end

heap:                      ; 堆内存的开始位置
    .space HEAP_SIZE       ; 开辟一块长度为 4096 字节的内存空间
heap_end:                  ; 堆内存的结束位置
\end{lstlisting}

这里的 \texttt{.space} 是汇编器的伪指令：它在目标文件里预留空间，
不产生任何机器码。运行时这就是进程虚拟地址空间中的一块内存。

注意我们还定义了两个缓冲区：

\begin{lstlisting}[style=asm,caption=hello.s 中的 I/O 缓冲区]
.global input_buf
input_buf: .space 256       ; 存放键盘读入的内容

.global output_buf
output_buf: .space 256      ; 存放准备输出的内容
\end{lstlisting}

这两块 256 字节的内存，分别对应命令行输入和程序输出。

\subsection{memory\_init：初始化堆}

分配器初始化只做一件事：在堆的第一个 8 字节头部填入整个堆的大小和"空闲"标记。

\begin{lstlisting}[style=asm,caption=memory.s — memory\_init]
.include "src/memory.inc"
.extern heap
.extern heap_end
.global memory_init
.global memory_alloc

memory_init:
    adrp x0, heap@PAGE           ; 拿到 heap 所在页的基地址
    add  x0, x0, heap@PAGEOFF    ; 加页内偏移，得到 heap 真正的虚拟地址
    mov  x1, #HEAP_SIZE          ; x1 = 堆总大小 = 4096
    sub  x1, x1, #8              ; x1 = 可用空间 = 4096 - 8 = 4088
    str  w1, [x0]                ; 把可用大小写入 head 的第一个字（size 字段）
    mov  w2, #0
    str  w2, [x0, #4]            ; 第二字（used 字段）写 0，表示空闲
    ret
\end{lstlisting}

执行完 \texttt{memory\_init}，堆的状态：

\begin{center}
\begin{tabular}{|c|c|c|}
\hline
4088 & 0 & (未初始化，4088 字节空闲) \\
\hline
\multicolumn{2}{|c|}{第一个块的头部} & \\
\hline
\end{tabular}
\end{center}

\paragraph{地址计算的关键：adrp + add}
ARM 64 位下不能用一条指令直接写一个 64 位绝对地址。
\texttt{adrp} 把 \texttt{heap} 所在的 4 KB 页的基地址写入 \texttt{x0}，
然后 \texttt{add} 再加上页内偏移——两步合起来就得到了真正的虚拟地址。
这是位置无关代码（PIC）的标准做法。

\subsection{memory\_alloc：分配一块内存}

这是分配器的核心。它的思路是顺序扫描（first-fit）：

\begin{lstlisting}[style=asm,caption=memory.s — memory\_alloc（第一部分：扫描循环）]
memory_alloc:
    ; 入栈保护：保存 x19-x22（调用者保存寄存器）
    stp x19, x20, [sp, #-32]!
    stp x21, x22, [sp, #16]
    mov x19, x0                   ; x19 = 申请的大小
    adrp x20, heap@PAGE
    add  x20, x20, heap@PAGEOFF   ; x20 = 扫描指针（从 heap 头部开始）
    adrp x21, heap_end@PAGE
    add  x21, x21, heap_end@PAGEOFF

alloc_loop:
    cmp  x20, x21                 ; 超过堆尾？
    b.hs alloc_fail               ; 是 → 失败

    ldr  w22, [x20]               ; 读当前块的 size
    ldr  w23, [x20, #4]           ; 读当前块的 used 标记

    cmp  w23, #1                  ; 当前块在用吗？
    b.eq next_block               ; 是 → 跳过
    cmp  w22, w19                 ; 当前块大小 >= 需要的大小？
    b.lt next_block               ; 不够 → 跳过
\end{lstlisting}

一旦找到一块合适的空闲块，我们需要决定是否把它\textbf{切割}成两块——
一块分配出去，剩下的留作新的空闲块。只有剩余空间足够容纳
一个头部（8 字节）+ 至少 1 字节数据时才有切割的意义。

\begin{lstlisting}[style=asm,caption=memory.s — memory\_alloc（第二部分：分配与切割）]
    ; 找到了：先标记为已用
    mov  w0, #1
    str  w0, [x20, #4]            ; used = 1

    sub  x1, x22, x19             ; x1 = 剩余 = 当前块大小 - 要分配的大小
    cmp  x1, #8                   ; 剩余小于 8 字节？
    b.le no_split                 ; 是 → 不切割，整块都分配

    ; 切割：在剩余位置写入一个新空闲块的头部
    add  x2, x20, #8
    add  x2, x2, x19              ; x2 = 新块头部 = 当前块头 + 8 + 分配大小
    sub  x3, x1, #8               ; x3 = 新块可用字节 = 剩余 - 8 (新块头)
    str  w3, [x2]                 ; 新块 size = x3
    mov  w4, #0
    str  w4, [x2, #4]             ; 新块 used = 0（空闲）

    str  w19, [x20]               ; 更新当前块 size = 申请的大小

no_split:
    add  x0, x20, #8              ; 返回值 = 数据区首地址（头 + 8）
    b    alloc_done

next_block:
    add  x20, x20, #8             ; 跳过头部
    add  x20, x20, x22            ; 跳过数据区，x20 指向了下一个块
    b    alloc_loop

alloc_fail:
    mov  x0, #0                   ; 返回 0（NULL），表示分配失败

alloc_done:
    ldp  x21, x22, [sp, #16]
    ldp  x19, x20, [sp], #32      ; 恢复寄存器，同时释放栈帧
    ret
\end{lstlisting}

\paragraph{为什么要保护 x19--x22？}
AArch64 的调用约定（AAPCS64）把寄存器分成两组：
\begin{itemize}
    \item \textbf{临时寄存器} \texttt{x0--x18}：函数可以随便用，
        用完不需要恢复。调用者如果还要用，必须自己保存。
    \item \textbf{被调用者保存寄存器} \texttt{x19--x28}：
        函数如果要用它们，必须在函数返回前恢复原值。
\end{itemize}

\texttt{memory\_alloc} 用了 \texttt{x19--x22} 来保存扫描指针和大小，
所以进入函数时先把它们压栈，退出前恢复。这是标准的函数序言/尾声。

\subsection{memory\_free：释放内存并合并相邻空闲块}

释放的逻辑是：把块标记为空闲，然后检查左右邻居是否也是空闲，
如果是就把它们合并成一块（否则堆会被切成越来越碎的小块）。

\begin{lstlisting}[style=asm,caption=memoryfree.s — memory\_free]
.include "src/memory.inc"
.extern heap
.extern heap_end
.global memory_free

memory_free:
    stp x19, x20, [sp, #-96]!
    stp x21, x22, [sp, #16]
    stp x23, x24, [sp, #32]
    stp x25, x26, [sp, #48]
    stp x27, x28, [sp, #64]
    stp x29, x30, [sp, #80]

    ; 第一步：标记当前块为空闲
    cbz x0, free_done             ; NULL 指针直接返回
    sub x19, x0, #8               ; x19 = 当前块头部
    mov w1, #0
    str w1, [x19, #4]             ; used = 0

    ; 准备扫描指针，向前一块找
    adrp x20, heap@PAGE
    add  x20, x20, heap@PAGEOFF
    adrp x21, heap_end@PAGE
    add  x21, x21, heap_end@PAGEOFF
    mov  x22, x20
    mov  x23, #0

find_prev_loop:
    cmp x22, x19                  ; 扫描到了当前块位置？
    b.hs check_prev               ; 是 → 跳出循环，此时 x23 就是前一块
    cmp x22, x21                  ; 超过堆尾？
    b.hs check_prev
    mov x23, x22                  ; 记录"上一块"的位置
    ldr w24, [x22]
    add x22, x22, #8
    add x22, x22, x24             ; 跳到下一个块的头
    b find_prev_loop
\end{lstlisting}

循环结束后，\texttt{x23} 里存的是前一块的地址（如果当前块是第一个块，
\texttt{x23} 保持为 0）。接下来检查前后两个邻居：

\begin{lstlisting}[style=asm,caption=memoryfree.s — 合并相邻块]
check_prev:
    cbz x23, check_next           ; 没有前一块 → 跳过
    ldr w25, [x23, #4]
    cmp w25, #0                   ; 前一块是空闲吗？
    b.ne check_next               ; 不空闲 → 跳过
    ; 合并前一块：把前一块的 size 扩大，把当前块吃掉
    ldr w24, [x23]                ; 前一块 size
    ldr w26, [x19]                ; 当前块 size
    add w24, w24, w26
    add w24, w24, #8              ; +8 吃掉当前块头部
    str w24, [x23]                ; 写回前一块的 size
    mov x19, x23                  ; 现在"当前块"就是合并后的那块

check_next:
    ldr w26, [x19]
    add x27, x19, #8
    add x27, x27, w26             ; x27 = 当前块后面那一块的头
    cmp x27, x21                  ; 后一块超过堆尾？
    b.hs free_done                ; 是 → 没东西合并

    ldr w28, [x27, #4]
    cmp w28, #0                   ; 后一块空闲吗？
    b.ne free_done
    ; 合并后一块：把后一块的 size 加到当前块
    ldr w29, [x27]
    add w26, w26, w29
    add w26, w26, #8
    str w26, [x19]

free_done:
    ; 恢复所有保存过的寄存器（与进入时的 stp 顺序相反）
    ldp x29, x30, [sp, #80]
    ldp x27, x28, [sp, #64]
    ldp x25, x26, [sp, #48]
    ldp x23, x24, [sp, #32]
    ldp x21, x22, [sp, #16]
    ldp x19, x20, [sp], #96
    ret
\end{lstlisting}

释放一个块时，最多可能合并三块（前一块 + 当前块 + 后一块）。
正因为我们扫描的方式——从堆头一路数块大小、逐个跳过——才能正确找到
"当前块的前一块"的头部位置。

\subsection{测试内存分配器}

\texttt{test.s} 里有一个专门测试内存系统的函数：

\begin{lstlisting}[style=asm,caption=test.s — test\_memory]
test_memory:
    stp x29, x30, [sp, #-16]!
    sub sp, sp, #32               ; 留出 32 字节栈空间存指针

    mov x0, #8
    bl memory_alloc               ; 第一次分配 8 字节
    str x0, [sp, #0]              ; 存返回的指针到栈上
    ldr w4, =0x11111111
    ldr w5, =0x11111111
    str w4, [x0]                  ; 写入测试数据
    str w5, [x0, #4]

    mov x0, #8
    bl memory_alloc               ; 第二次分配（验证切割是否正常）
    str x0, [sp, #8]
    ldr w4, =0x22222222
    ...

    mov x0, #8
    bl memory_alloc               ; 第三次分配
    str x0, [sp, #16]
    ldr w4, =0x33333333
    ...

    ldr x0, [sp, #16]
    bl memory_free                ; 先释放第三块
    str xzr, [sp, #16]

    ldr x0, [sp, #8]
    bl memory_free                ; 再释放第二块（触发与第三块合并）
    str xzr, [sp, #8]

    add sp, sp, #32
    ldp x29, x30, [sp], #16
    ret
\end{lstlisting}

这个测试的意义是：如果分配、切割、释放、合并这四个步骤都正确工作，
堆里的状态就会经历"三块在用"→"第一块在用，后两块释放并合并"的过程。
后续再分配一个大请求时，就会从合并后的大块里切分——这正是分配器的正常工作方式。
